\documentclass{article}
\usepackage[UTF8]{ctex}
\usepackage{amsmath,amssymb,amsfonts}
\usepackage{enumitem}
\usepackage{geometry}
\geometry{a4paper, margin=1in}
% === 错题本排版优化（不参与题目标识） ===
\setlength{\parindent}{0pt}
\setlength{\parskip}{0.35em plus 0.10em minus 0.08em}
\setlength{\emergencystretch}{3em}
\setlist{itemsep=0.25em, topsep=0.40em, parsep=0pt}
\allowdisplaybreaks[2]
% === 错题本排版优化结束 ===
\begin{document}
\section*{上海大学2025年数学分析真题}
\section*{46．上海大学2025年数学分析试卷}

\subsection{题目1}
1．求数列极限：（1） \(\lim _{n \rightarrow \infty} \sqrt[n]{n}\) ，（2） \(\lim _{n \rightarrow \infty}(2 n!)^{\frac{1}{n^{2}}}\) 。

\subsection{题目2}
2．设函数 \(f(x)\) 在 \([-1,1]\) 上连续，求 \(\lim _{x \rightarrow 0} \frac{\ln (1+f(x) \arctan x)}{2^{x}-1}\) ．

\subsection{题目3}
3．求曲线积分

\[
\oint_{L} \frac{x^{2} y}{(x-a)^{2}+y^{2}} \mathrm{~d} y
\]

其中 \(L\) 是由圆 \((x-a)^{2}+y^{2}=a^{2}(a>0)\) 与 \(x\) 轴围成的第四象限区域的边界，方向为逆时针。

\subsection{题目4}
4．计算曲面积分

\[
\iint_{\sum} x^{2} \mathrm{~d} y \mathrm{~d} z+y \mathrm{~d} z \mathrm{~d} x+z \mathrm{~d} x \mathrm{~d} y .
\]

其中 \(\sum\) 为曲面 \(z=\sqrt{x^{2}+y^{2}}\) 在 \(0 \leq z \leq 1\) 的部分，方向取下侧。

\subsection{题目5}
1．若函数 \(f(x)\) 为区间 \([a, b]\) 上只有有限个间断点的有界函数，证明：\(f(x)\) 在 \([a, b]\) 上可积．

\subsection{题目6}
2．设函数 \(f(x)\) 在 \([0,1]\) 上二阶可微，\(f(0)=f(1)=0, \min _{0 \leq x \leq 1} f(x)=-1\) ．证明：

\[
\min _{0 \leq x \leq 1} f^{\prime \prime}(x) \leq 8, \max _{0 \leq x \leq 1} f^{\prime \prime}(x) \geq 8 .
\]

\subsection{题目7}
3．对于级数 \(\sum_{n=1}^{\infty}(-1)^{n} \frac{x^{2}+2^{n}}{(n+1)!}\) 。

\begin{itemize}
  \item[（1）] 证明：在任意有限区间 \(I\) 上一致收敛。
  \item[（2）] 判断级数在任何一点的敛散性，若收敛，说明是条件收敛还是绝对收敛，并给出理由。
\end{itemize}

\subsection{题目8}
4．解答如下问题：

\begin{itemize}
  \item[（1）] 叙述 \(\mathbb{R}^{2}\) 的聚点定理．
  \item[（2）] 利用聚点定理证明二元函数 \(f\) 在有界闭域 \(D \subset \mathbb{R}^{2}\) 连续则一致连续．
\end{itemize}

\subsection{题目9}
5．设函数 \(\varphi(x)\) 在 \([a, b]\) 上有界可积，证明： \(\lim _{p \rightarrow+\infty} \int_{a}^{b} \varphi(x) \sin p x \mathrm{~d} x=0\) ．

\subsection{题目10}
6．设函数 \(f(x, y)= \begin{cases}\frac{x y^{2}}{\sqrt{x^{2}+y^{2}}}, & (x, y) \neq(0,0) ; \\ 0, & (x, y)=(0,0) .\end{cases}\)

\begin{itemize}
  \item[（1）] 计算 \(f(x, y)\) 的偏导数．
  \item[（2）] 讨论 \(f(x, y)\) 在原点的连续性．
  \item[（3）] 讨论 \(f(x, y)\) 在原点的可微性．
\end{itemize}

\subsection{题目11}
7．解答如下问题：\\
（1）若函数 \(f(x)\) 满足以下条件：（i）函数 \(f(x)\) 在点 \(x_{0}\) 的邻域 \(U\left(x_{0}\right)\) 连续；（ii）在 \(U^{\circ}\left(x_{0}\right)\) 可导；（iii）极限 \(\lim _{x \rightarrow x_{0}} f^{\prime}(x)\) 存在．证明：\(f(x)\) 在点 \(x_{0}\) 处可导，且 \(f^{\prime}\left(x_{0}\right)=\lim _{x \rightarrow x_{0}} f^{\prime}(x)\) ．\\
（2）若函数 \(f(x)\) 只满足本题（1）中条件（i）和（ii），但不满足（iii），则 \(f(x)\) 在点 \(x_{0}\) 不可导，是否正确？请给出理由。
\end{document}
